\subsubsection{Period $N = 60$: the $5 \times 5$ pentagonal family}

The longest period in the classification is $N = 60$, occurring on
$52$ parameter sets across the $5 \times 5$, $5 \times 10$, and
$10 \times 10$ tori. The reflection parameters are the golden-ratio
surs $(5-\sqrt{5})/10$ and $(5+\sqrt{5})/10$, and the phase parameters
take values in $\{0, 2\pi/5, 4\pi/5\}$. The $5 \times 5$ torus alone
contains $36$ of these parameter sets, making it by far the richest
single torus in the classification. The value
$(5-\sqrt{5})/10 \approx 0.2764$ admits an elegant characterisation in
terms of the golden ratio $\varphi = (1+\sqrt{5})/2$, since it satisfies
$\rho = 1/(\varphi + 2)$. This is the clearest instance in the data of
an inverse relationship between torus size and maximum revival period,
and it is the single most striking result of the classification.
